\clearpage
\section{光线与场景中的三角形}
雾里的一点能否被灯照亮，取决于它和灯之间有没有墙。场景里有几十万个三角形，逐个检查会很慢。加载模型时，CPU 给一组相邻的三角形套上长方体，称为包围盒；再把这组分小，直到一个盒子里最多剩下六个三角形。这样组成的树叫作包围盒层次结构（BVH）。GPU 查遮挡时，先试光线能否穿过大盒子。连盒子都碰不到，里面的三角形就可以一起跳过。

\begin{center}
\begin{tikzpicture}[x=.92cm,y=.62cm,font=\small]
\foreach \y/\labeltext in {2.4/各取八个,0/沿空隙分开}{
  \node[anchor=east] at (-.35,\y+.5) {\labeltext};
  \foreach \x in {0,1,2,9}{
    \fill[quiet!15] (\x,\y+.07) rectangle ++(1,.86);
    \draw[quiet!60] (\x,\y+.07) rectangle ++(1,.86);
    \draw[quiet!60] (\x,\y+.07) -- ++(1,.86) (\x,\y+.93) -- ++(1,-.86);
  }
}
\draw[wine,thick,rounded corners=1pt] (-.04,2.32) rectangle (1.97,3.48);
\draw[teal!70!black,thick,rounded corners=1pt] (2.02,2.32) rectangle (10.04,3.48);
\draw[wine,thick,rounded corners=1pt] (-.04,-.08) rectangle (3.04,1.08);
\draw[teal!70!black,thick,rounded corners=1pt] (8.96,-.08) rectangle (10.04,1.08);
\foreach \y in {0,2.4}{\draw[->,thick] (5.5,\y+1.6) -- (5.5,\y+.35);}
\end{tikzpicture}
\end{center}
{\small\color{quiet}俯视示意，按本项目的划分方法绘制。每个灰格代表四个三角形，彩框是分出来的两个盒子。箭头穿过同一条空走廊。上图还要进入右边的盒子查找，下图在两个盒外就能排除遮挡。}\par

怎样估算一种划分要做多少工作？从许多方向发出射线时，较大的盒子通常更容易被穿过。设左右盒子的表面积为 $S_L,S_R$，其中的三角形数为 $N_L,N_R$，便可用 $S_LN_L+S_RN_R$ 比较同一个父盒里的划分。图中每个灰格的长、宽、高都是 1。上图的盒长分别为 2 和 8，表面积分别为 10 和 34，得分为 $10\times8+34\times8=352$。下图的盒长为 3 和 1，得分降为 $14\times12+6\times4=192$。两图总共有同样多的三角形，差别来自盒子包进了多少空地。

在 \texttt{GIScene.cpp} 的 \texttt{SplitTriangles} 中，CPU 沿三角形中心分布最宽的方向分成 12 段，比较中间的 11 个分界。每比较一次只需要两侧的盒子与个数。代码用半个表面积乘个数，所有得分同时除以二，选出来的位置相同。三角形很少或中心重合时，直接平分。画面渲染期间，GPU 沿建好的树查找。查阴影只要在灯前遇到一块不透明表面就能停，查反射则要保留最近交点，才能读到正确的材质。

{\small\color{quiet}\href{https://pbr-book.org/4ed/Primitives_and_Intersection_Acceleration/Bounding_Volume_Hierarchies}{PBRT 第四版 §7.3.2，按表面积划分包围盒}
\quad\href{run:Redflag-Debugging-Notes.pdf}{墙面接缝的排错图}
\quad\href{run:Redflag-Rendering-Performance.pdf}{两座建筑的 GPU 实测}}\par
